\documentclass[lang=cn,newtx,10pt,scheme=chinese]{elegantbook}


\begin{document}

\tikzset{
  point/.style = {% define a style for the function points
    circle,
    fill=#1,
    draw=black,
    inner sep=1pt,
  },
  point/.default = {violet!60}
}

\begin{figure}[!ht]
\centering
\resizebox{0.4\textwidth}{!}{%
\begin{tikzpicture}[declare function={f(\x)=(\x-5)^2;}]
    \draw[thick, dashed](0,0)--(15,0) node[right]{$\theta$};
    \draw[thick, dashed](0,0)--(0,30) node[above]{Cost};
    \draw[thick, dashed](0.6, {f(0.6)}) -- (0.6,0)
          node[below, text width=4em, align=center]{Random initial value};
    \draw[thick, dashed](3, {f(3)}) -- (3,0) node[below]{$\hat\theta$};
    \draw[domain=4:6, smooth, thick, gray] plot (\x, f(\x);
    \foreach \a [count=\c, remember=\c as \C] in {4, 4.2, 4.36, 4.488, 4.5904} {
      \node[point] (\c) at (\a, {f(\a)}){};
      \ifnum\c>1 % after the first coordinate draw an arrow
         \draw[->, bend left](\C) to (\c);
      \fi
    }
 \end{tikzpicture}
}
\end{figure}

\end{document}